%% ----------------------------------------------------------------------------
%%  §9.3  L'integrale indefinito. Integrali immediati. Decomposizione in somma
%% ----------------------------------------------------------------------------
\section{L'integrale indefinito}
\label{sec:integrale-indefinito}

Nel paragrafo precedente abbiamo ricondotto il calcolo di un integrale definito alla
ricerca di una primitiva della funzione integranda. È perciò naturale porre la
definizione seguente.

\begin{definizione}[Integrale indefinito]\label{def:integrale-indefinito}
\index{integrale indefinito}
Sia $f$ una funzione continua in un intervallo $I$. L'insieme di tutte le primitive di
$f$ in $I$ si chiama \emph{integrale indefinito} di $f$ e si indica con il simbolo
\begin{equation}\label{eq:integrale-indefinito}
  \int f(x)\,dx .
\end{equation}
\end{definizione}

Per il \cref{teo:caratterizzazione-primitive}, se $F$ è una primitiva di $f$ in $I$,
\begin{equation}\label{eq:integrale-indefinito-c}
  \int f(x)\,dx=F(x)+c ,
\end{equation}
dove $c$ è una costante arbitraria: la scrittura significa che l'insieme delle
primitive è $\set{F+c \tc c\in\R}$. C'è una differenza sostanziale tra l'integrale
definito $\int_{a}^{b}f(x)\,dx$, che è un \emph{numero}, e l'integrale indefinito $\int f(x)\,dx$,
che è un \emph{insieme di funzioni}; il legame tra i due è la formula fondamentale
\eqref{eq:formula-fondamentale}. Nella pratica si scrive $\int f(x)\,dx=F(x)+c$ e si dice,
con un abuso di linguaggio, che $F$ «è» l'integrale indefinito di $f$.

Poiché la derivata di una somma è la somma delle derivate e la derivata di $cF$ è
$cF'$ (\cref{teo:operazioni-derivate}), se $F$ e $G$ sono primitive di $f$ e $g$, allora
$F+G$ è una primitiva di $f+g$ e $cF$ è una primitiva di $cf$. In termini di integrali
indefiniti,
\begin{equation}\label{eq:linearita-indefinito}
\begin{gathered}
  \int\bigl[f(x)+g(x)\bigr]dx=\int f(x)\,dx+\int g(x)\,dx,\\
  \int c\,f(x)\,dx=c\int f(x)\,dx \quad (c\neq0),
\end{gathered}
\end{equation}
in analogia con la linearità dell'integrale definito (\cref{teo:linearita}). Le
uguaglianze vanno intese a meno di costanti additive.

\subsection{Integrali immediati}

Leggendo da destra a sinistra la tabella delle derivate elementari
(\cref{tab:derivate-elementari}) e le derivate delle funzioni trigonometriche inverse
\eqref{eq:derivata-arcsen}, \eqref{eq:derivata-arctg} si ottengono gli \emph{integrali immediati} della
\cref{tab:integrali-immediati}. Ciascuno si verifica derivando il secondo membro, e
ciascuno vale in ogni intervallo in cui la funzione integranda è continua.

\begin{table}[htb]
\centering
\renewcommand{\arraystretch}{1.9}
\begin{tabular}{@{}l@{\qquad}l@{}}
\toprule
$\displaystyle\int x^{b}\,dx=\frac{x^{b+1}}{b+1}+c \quad (b\neq-1)$ &
$\displaystyle\int\frac{dx}x=\log\abs x+c$\\
$\displaystyle\int e^{x}\,dx=e^{x}+c$ &
$\displaystyle\int a^{x}\,dx=\frac{a^{x}}{\log a}+c \quad (a>0,\ a\neq1)$\\
$\displaystyle\int\sen x\,dx=-\cos x+c$ &
$\displaystyle\int\cos x\,dx=\sen x+c$\\
$\displaystyle\int\frac{dx}{\cos^{2}x}=\tg x+c$ &
$\displaystyle\int\frac{dx}{\sqrt{1-x^{2}}}=\arcsen x+c$\\
$\displaystyle\int\frac{dx}{1+x^{2}}=\operatorname{arctg}x+c$ & \\
\bottomrule
\end{tabular}
\caption{Integrali immediati. Ciascuna formula vale in ogni intervallo in cui la
funzione integranda è continua: ad esempio, la prima con $b=-2$ vale in $(-\infty,0)$ e
in $(0,+\infty)$, la seconda in ciascuno dei due intervalli che non contengono lo zero.}
\label{tab:integrali-immediati}
\end{table}

A proposito del logaritmo, osserviamo che
\begin{equation}\label{eq:derivata-log-abs}
  D\log\abs x=\frac1x \qquad\perogni x\neq0 .
\end{equation}
Per $x>0$ è la \eqref{eq:derivata-log}. Per $x<0$ si ha $\abs x=-x$, e per la regola di
derivazione delle funzioni composte
\[
  D\log\abs x=D\log(-x)=\frac1{-x}\cdot(-1)=\frac1x \qquad (x<0) .
\]
Dunque $\int\frac{dx}x=\log\abs x+c$, intendendo che l'integrale è considerato in un intervallo
che non contiene il punto $x=0$.

\subsection{Integrali immediati di funzioni composte}

Molti integrali si riconducono a quelli della tabella grazie alla regola di derivazione
delle funzioni composte (\cref{teo:derivata-composta}). Se $G$ è una primitiva di $g$ e
$\varphi$ è derivabile, allora $D\,G(\varphi(x))=g(\varphi(x))\,\varphi'(x)$, cioè
\begin{equation}\label{eq:composta-indefinito}
  \int g\bigl(\varphi(x)\bigr)\,\varphi'(x)\,dx=G\bigl(\varphi(x)\bigr)+c .
\end{equation}
In particolare, per una funzione $\varphi$ derivabile (positiva nella prima formula se $b$
non è intero, non nulla se $b$ è intero negativo; non nulla nella seconda),
\begin{equation}\label{eq:integrali-composti}
\begin{aligned}
  \int\varphi(x)^{b}\,\varphi'(x)\,dx&=\frac{\varphi(x)^{b+1}}{b+1}+c \qquad (b\neq-1),\\
  \int\frac{\varphi'(x)}{\varphi(x)}\,dx&=\log\abs{\varphi(x)}+c,\\
  \int e^{\varphi(x)}\varphi'(x)\,dx&=e^{\varphi(x)}+c,\\
  \int\frac{\varphi'(x)}{1+\varphi^{2}(x)}\,dx&=\operatorname{arctg}\varphi(x)+c .
\end{aligned}
\end{equation}
Il caso più frequente è quello di un argomento lineare, $\varphi(x)=\alpha x+\beta$ con
$\alpha\neq0$: allora $\varphi'=\alpha$, e
\begin{equation}\label{eq:argomento-lineare}
  \int g(\alpha x+\beta)\,dx=\frac1\alpha\,G(\alpha x+\beta)+c ;
\end{equation}
ad esempio, $\int\cos2x\,dx=\frac12\sen2x+c$ e $\int e^{-3x}\,dx=-\frac13e^{-3x}+c$.

\begin{esempio}[Integrali di funzioni composte]\label{es:integrali-composti}
\begin{enumerate}[label=\textup{(\alph*)}]
  \item In $\tg x=\frac{\sen x}{\cos x}$ il numeratore è, a meno del segno, la derivata
  del denominatore:
  \begin{equation}\label{eq:integrale-tg}
    \int\tg x\,dx=-\int\frac{-\sen x}{\cos x}\,dx=-\log\abs{\cos x}+c .
  \end{equation}
  \item $\int\frac x{1+x^{2}}\,dx=\frac12\int\frac{2x}{1+x^{2}}\,dx=\frac12\log(1+x^{2})+c$ (il valore
  assoluto non serve, perché $1+x^{2}>0$).
  \item $\int x\,e^{x^{2}}dx=\frac12\int e^{x^{2}}\,2x\,dx=\frac12e^{x^{2}}+c$.
  \item $\int\sen x\cos x\,dx=\int\sen x\,D(\sen x)\,dx=\frac12\sen^{2}x+c$.
  \item $\int\frac{dx}{x\log x}=\int\frac{1/x}{\log x}\,dx=\log\abs{\log x}+c$, in $(0,1)$ o in
  $(1,+\infty)$.
\end{enumerate}
\end{esempio}

\subsection{Integrazione per decomposizione in somma}

In molti casi il calcolo dell'integrale indefinito di una funzione si riconduce a quello
di integrali già noti, o più semplici, decomponendo la funzione integranda nella somma
di due o più funzioni e applicando la linearità \eqref{eq:linearita-indefinito}.

\begin{esempio}[Quattro decomposizioni]\label{es:decomposizione}
\begin{enumerate}[label=\textup{(\alph*)}]
  \item Sommando e sottraendo $1$ al numeratore:
  \[
  \begin{aligned}
    \int\frac x{x+1}\,dx&=\int\frac{x+1-1}{x+1}\,dx=\int\Bigl(1-\frac1{x+1}\Bigr)dx\\
    &=x-\log\abs{x+1}+c .
  \end{aligned}
  \]
  \item Poiché $\tg^{2}x=\frac{\sen^{2}x}{\cos^{2}x}=\frac{1-\cos^{2}x}{\cos^{2}x}=\frac1{\cos^{2}x}-1$,
  \[
    \int\tg^{2}x\,dx=\int\Bigl(\frac1{\cos^{2}x}-1\Bigr)dx=\tg x-x+c .
  \]
  \item Scrivendo $1=\sen^{2}x+\cos^{2}x$ al numeratore, la frazione si scinde:
  \[
  \begin{aligned}
    \int\frac{dx}{\sen x\cos x}&=\int\Bigl(\frac{\sen x}{\cos x}+\frac{\cos x}{\sen x}\Bigr)dx\\
    &=-\log\abs{\cos x}+\log\abs{\sen x}+c=\log\abs{\tg x}+c .
  \end{aligned}
  \]
  \item Dalla formula di duplicazione $\cos2x=\cos^{2}x-\sen^{2}x=1-2\sen^{2}x$
  \eqref{eq:duplicazione} si ha $\sen^{2}x=\frac{1-\cos2x}2$, e per la
  \eqref{eq:argomento-lineare}
  \begin{equation}\label{eq:integrale-sen2}
  \begin{split}
    \int\sen^{2}x\,dx&=\frac12\int(1-\cos2x)\,dx=\frac x2-\frac{\sen2x}4+c\\
    &=\frac12\,(x-\sen x\cos x)+c ,
  \end{split}
  \end{equation}
  dove nell'ultimo passaggio si è usata la duplicazione $\sen2x=2\sen x\cos x$. Ne segue
  \begin{equation}\label{eq:integrale-cos2}
    \int\cos^{2}x\,dx=\int(1-\sen^{2}x)\,dx=\frac12\,(x+\sen x\cos x)+c .
  \end{equation}
\end{enumerate}
\end{esempio}

Il calcolo simbolico permette di verificare rapidamente un integrale indefinito. Il
programma seguente usa \texttt{sympy}: per ciascuna funzione calcola una primitiva $F$
e controlla che $F'-f$ si semplifichi a zero.

\begin{python}[Integrali indefiniti con sympy, e la loro verifica]
import sympy as sp

x = sp.symbols("x", real=True)
esempi = [x / (x + 1), sp.tan(x)**2, sp.sin(x)**2, sp.tan(x),
          x * sp.exp(x**2), 1 / (x * sp.log(x))]
for f in esempi:
    F = sp.integrate(f, x)
    verifica = sp.simplify(sp.diff(F, x) - f)
    print(f"f = {f}")
    print(f"    F = {F}      F' - f = {verifica}")
\end{python}
\begin{risultato}
f = x/(x + 1)
    F = x - log(x + 1)      F' - f = 0
f = tan(x)**2
    F = -x + sin(x)/cos(x)      F' - f = 0
f = sin(x)**2
    F = x/2 - sin(x)*cos(x)/2      F' - f = 0
f = tan(x)
    F = -log(cos(x))      F' - f = 0
f = x*exp(x**2)
    F = exp(x**2)/2      F' - f = 0
f = 1/(x*log(x))
    F = log(log(x))      F' - f = 0
\end{risultato}

Si osservi che \texttt{sympy} non scrive la costante arbitraria e scrive $\log(x+1)$,
$\log(\cos x)$, $\log(\log x)$ senza valore assoluto: le sue primitive sono corrette
negli intervalli in cui l'argomento del logaritmo è positivo (per $\cos x<0$, ad
esempio, \texttt{sympy} lavora in campo complesso). Il risultato di un programma di
calcolo simbolico va sempre letto con senso critico, controllando l'intervallo in cui
vale.

\begin{notastorica}[Le tavole degli integrali e i limiti del calcolo]
Tra le prime raccolte sistematiche di integrali indefiniti vi furono quelle delle
\emph{Institutiones calculi integralis} di Eulero (1768--1770), che fissarono per più di
un secolo i metodi che studieremo nei prossimi paragrafi. Ci si accorse presto che
alcune funzioni elementari molto semplici, come $e^{-x^{2}}$, $\frac{\sen x}x$, $\frac1{\log x}$,
resistevano a ogni tentativo di integrazione: Joseph Liouville dimostrò, in una serie di
memorie del 1833--1841, che le loro primitive non si possono esprimere mediante un numero
finito di funzioni elementari. Le primitive esistono (per il teorema fondamentale), ma
definiscono funzioni nuove, come la funzione degli errori $\operatorname{erf}$, che si
calcolano numericamente. Nel 1969 Robert Risch pubblicò un algoritmo (completo nel caso delle funzioni
trascendenti) che decide se una funzione elementare ha una primitiva elementare e, in
caso affermativo, la calcola: i programmi di calcolo simbolico come \texttt{sympy} ne
implementano versioni parziali, affiancate da procedimenti euristici.
\end{notastorica}
